\subsection{qLMPC Optimization Algorithm}\label{sec:qLMPCOptimizationAlgo}
The predictive control law is formulated as an optimization problem, minimizing the finite horizon objective function \( J_k \) at iteration \( k \) over a predefined number of \( n_h \) time steps:
\begin{align*}
    J_k(\bm{e}_d, \Delta \bm{u}_{\text{Mdl}}) = & \underbrace{ \sum_{i=k}^{k+n_h-2}  \bm{e}_{d}^T(i) \bm{Q}_{\text{WT}} \bm{e}_{d}(i)}_{J_Q(\Delta \bm{u}_{\text{Mdl}})} + \underbrace{\sum_{i=k}^{k+n_h-1} \Delta \bm{u}^T_{\text{Mdl}}(i) \bm{R}_{\text{WT}} \Delta \bm{u}_{\text{Mdl}}(i)}_{J_R(\Delta \bm{u}_{\text{Mdl}})} \\
    &+ \underbrace{\bm{e}_{d}(k+n_h-1)^T \bm{P}_{\text{WT}} \bm{e}_{d}(k+n_h-1)}_{J_P(\Delta \bm{u}_{\text{Mdl}})},
\end{align*}
where \( \bm{e}_{d} \in \mathbb{R}^{n_x \times 1} \) represents the error between the augmented, discrete state \( \bm{x}_{d}\) and its steady state vector $\bm{x}_{s}$ at step \( i \),
\begin{align*}
    \bm{e}_{d}(i) = \bm{x}_{d}(i) - \bm{x}_{s}(i),
\end{align*}
while $\bm{\Delta u}_{\text{Mdl}}(i)$ denotes the change in input  $\bm{u}_{\text{Mdl}}$ compared to the previous time step:
\begin{align*}
    \bm{\Delta u}_{\text{Mdl}}(i) = \bm{u}_{\text{Mdl}}(i) - \bm{u}_{\text{Mdl}}(i-1).
\end{align*}
The matrices \( \bm{Q}_{\text{WT}} \in \mathbb{R}^{n_x \times n_x} \) and \( \bm{R}_{\text{WT}} \in \mathbb{R}^{n_u \times n_u} \) are weighting matrices for the state errors and the input changes over the prediction horizon. The matrix \( \bm{P}_{\text{WT}} \in \mathbb{R}^{n_x \times n_x} \) penalizes the state error at the end of the trajectory. In this section, the dimensions are explicitly specified for the qLMPC design with \( n_x = 8 \) and \( n_u = 2 \).

The cost \( J_k \) consists of two parts: the stage cost $J_Q + J_R$, which represents the weighted sum of squares of the errors and inputs over the trajectory length, and the terminal cost $J_P$, which is the weighted sum of squares of the terminal states at sample $k+n_h-1$.
The constrained optimization problem to be solved in each sample can be formulated based on this objective function as
\begin{align} \label{eq:ObjFunWT}
    \begin{array}{ll}
        \underset{\Delta \bm{u}_{\text{Mdl}}}{\text{minimize }} & J_k(\bm{x}_d, \bm{x}_{s}, \Delta \bm{u}_{\text{Mdl}}) \\[10pt]
        \text{subject to} & \bm{x}_d(k+1) = \bm{A}_{d,k}(\bm{\rho}_{P,k}) \bm{x}_d(k) + \bm{B}_{d,k}(\bm{\rho}_{P,k}) \Delta \bm{u}_{\text{Mdl}}(k) \\[10pt]
        & \bm{e}_d(k) = \bm{C}_d \bm{x}_d(k) -\bm{x}_s(k).
    \end{array}
\end{align}
with $\bm{C}_d = \bm{I}_{8}$ and the steady-state solution from Equation~\eqref{eq:steadystatevec} with the first two elements of the vector $\bm{x}_s$ corresponding to the desired generator torque and blade pitch and the other six elements set to zero.
The weighting matrices \( \bm{Q}_{\text{WT}} \), \( \bm{R}_{\text{WT}} \), and \( \bm{P}_{\text{WT}} \) are diagonal matrices with the following values on their diagonals to penalize the state vector in Equation~\eqref{eq:AugmentedStateVector}:
\begin{align}
    \bm{q}_{\text{WT,diag}} &= [1 \quad 10000 \quad 0 \quad 1000 \quad 1000 \quad 0 \quad 0 \quad 0]^T \in \mathbb{R}^{8 \times 1}, \\
    \bm{r}_{\text{WT,diag}} &= [1 \quad 10000]^T \in \mathbb{R}^{2 \times 1}, \\
    \bm{p}_{\text{WT,diag}} &= 1000 \; \bm{q}_{\text{WT,diag}} \in \mathbb{R}^{8 \times 1}.
\end{align}
These matrices are designed to accommodate the different orders of magnitude of the physical states as well as their relationship among each other. The first two elements of the vector $\bm{q}_{\text{WT,diag}}$ are selected according to the order of magnitude between the generator torque reference input and the pitch angle reference input. The zeros in the weighting vectors $\bm{q}_{\text{WT,diag}}$ and $\bm{p}_{\text{WT,diag}}$ correspond to the slip angle and the three changes in input: generator torque, blade pitch, and wind speed. Penalizing the slip angle change, the third element in the state vector, is unnecessary since changes in rotor and generator speeds are penalized with the fourth and fifth element in the vector $\bm{q}_{\text{WT,diag}}$. These effectively capture the dynamics of the slip angle itself. Additionally, the weights on the control inputs are set to zero as changes in the control inputs are already penalized with the weighting vector $\bm{r}_{\text{WT,diag}}$. The weight on changes in the wind speed is set to zero because the wind is assumed to be constant during the control horizon. Moreover, it cannot be influenced by the control system and therefore there is no need to penalize changes in this variable. 

At each time step, the scheduling trajectory 
\begin{align}
    \bm{\Rho}_{k}= [\bm{\rho}_{\Rho}(k) \quad \bm{\rho}_{\Rho}(k+1)\quad \cdots \quad \bm{\rho}_{\Rho} (k+n_h-1)] \in \mathbb{R}^{3n_h \times 1}
\end{align}
is calculated based on the current scheduling parameter vector
\begin{align*}
    \bm{\rho}_{\Rho,k} = [\omega_r(k)\quad \beta_0(k) \quad V_\infty(k)] \in \mathbb{R}^{3 \times 1}
\end{align*}
and the estimated future values of the rotor speed, the pitch angle and the wind speed. For the controller initialization, the current values of the scheduling parameters are kept constant for all parameters in the trajectory, resulting in $\bm{\Rho}_0 = \bm{1}_{n^h} \otimes \bm{\rho}_{\Rho,k}$. Based on this 'frozen' scheduling parameter trajectory, an optimal control input trajectory can be calculated for Equation~\eqref{eq:ObjFunWT}. With this optimized control input trajectory, a new scheduling parameter trajectory $\bm{\Rho}_0^1$ is calculated and stored. This could be repeated until convergence is achieved, i.e., $\bm{\Rho}_0^{l-1} = \bm{\Rho}_0^l$ at the $l^\text{th}$ iteration step of the repetition of the optimization calculation. In practice, the number of iteration is predefined to guarantee the availability of a control signal within the controller sample time. For all following iterations, the estimated trajectory from the last controller iteration is used to calculate the control input trajectory of the current trajectory, that is the initial 'frozen' trajectory at time step $k$ is $\bm{\Rho}_{k-1}^{l} = \bm{\Rho}_k^0$. In \cite{hespe2021convergence} the convergence to an optimal trajectory is discussed: It is shown that guaranteeing convergence comes at the price of a high number of iterations and hence longer computation times. However, the suboptimal solutions that are reached after a few iterations are shown to be satisfactory for practical implementations, as is the case in this study.

The matrices \( \bm{\Lambda} \in \mathbb{R}^{8n_h \times 8}\) and \( \bm{S}  \in \mathbb{R}^{8n_h \times 2n_h} \) needed to calculate the optimal future control inputs are constructed by multiplication and concatenation of the system and input matrices from Equation~\eqref{eq:AugmentedStateSpace} as follows: % %\label{eq:AugmentedStateSpace}
%based on the augmented state vector, as provided in Equation~\eqref{eq:AugmentedStateVector}, and inputs, see Equation~\eqref{eq:DeltaUinput},
\begin{equation} \label{eq:nearlyToeplitzMat}
    \begin{aligned}
        \bm{\Lambda}(\bm{\Rho}_k) &= \begin{bmatrix}
            \bm{A}_d(\bm{\rho}_{P,k})  \\
            \bm{A}_d(\bm{\rho}_{P,k+1}) \bm{A}_d(\bm{\rho}_{P,k})  \\
            \vdots\\
            \bm{A}_d(\bm{\rho}_{P,k+n_h-1}) \bm{A}_d(\bm{\rho}_{P,k+n_h-2}) \cdots \bm{A}_d(\bm{\rho}_{P,k}) 
        \end{bmatrix}, \\[10pt]
        \bm{S}(\bm{\Rho}_k) &= 
        \begin{bmatrix}
            \bm{B}_d(\bm{\rho}_{P,k})  & \cdots & 0 \\
            \bm{A}_d(\bm{\rho}_{P,k+1}) \bm{B}_d(\bm{\rho}_{P,k}) & \cdots & 0 \\
            \vdots &  \ddots & \vdots \\
            \bm{A}_d(\bm{\rho}_{P,k+n_h-1}) \cdots \bm{A}_d(\bm{\rho}_{P,k+1}) \bm{B}_d(\bm{\rho}_{P,k}) & \cdots & \bm{B}_d(\bm{\rho}_{P,k+n_h-1})
        \end{bmatrix}.
    \end{aligned}
\end{equation}
The components of the cost function that take into account all $n_h$ samples of possible trajectories for the time steps $k$ to $k + n_h-1$ can be reformulated in vector formulation and calculated with the matrices from Equation~\eqref{eq:nearlyToeplitzMat} 

\begin{align*}
    J_{QP}(\Delta U) &= \left(\bm{\Lambda}(\bm{\Rho}_k)x_d(k) + \bm{S}(\bm{\Rho}_k) \Delta U(k)  
    - X_{s}\right)^T \bm{\mathcal{Q}}_{\text{WT}}  \left(\bm{\Lambda}(\bm{\Rho}_k)x_d(k) + \bm{S}(\bm{\Rho}_k) \Delta U(k) - X_{s}\right)\\
    J_R(\Delta U) &= \Delta U(k)^T \bm{\mathcal{R}}_{\text{WT}} \Delta U(k)
\end{align*}
where the penalty weight on the states is $J_{QP} = J_Q + J_P  \in \mathbb{R}^{8 \times 8}$ and with the matrices
\begin{align*}
    \bm{\mathcal{Q}}_{\text{WT}} =\begin{bmatrix} \bm{I}_{n_h -1} \otimes \bm{Q}_{\text{WT}} & \bm{0}_{8\times8(n_h-1)} \\ \bm{0}_{8(n_h-1)\times8}  & \bm{P}_{\text{WT}}\end{bmatrix} \in \mathbb{R}^{8n_h \times 8n_h}\quad\text{and}\quad   \bm{\mathcal{R}}_{\text{WT}} = \bm{I}_{n_h} \otimes \bm{R}_{\text{WT}}.
\end{align*}
The reference state is constant over the whole trajectory with $X_s = \bm{1}_{n_h-1 \times 1 }\otimes \bm{x}_s$, and the future control trajectory is defined as
\begin{align*}
    \Delta U(k) = \left[\Delta u(k)^T \quad \Delta u(k+1)^T \quad   \cdots  \quad   \Delta u(k+ n_h -1)^T \right]^T \in \mathbb{R}^{2n_h \times 1}.
\end{align*}

The initial state of the cost function, the state $\bm{x}_d$ from Equation~\eqref{eq:AugmentedStateVector} at sample $k$, can be calculated from the known state and control input at the previous sample $k-1$. The following property is used to derive the optimal control trajectory $\Delta U^*$ via solving a QP problem: For the constrained optimization problem with an objective function $V: \mathbb{R}^{n_{u} \times 1} \to \mathbb{R}$ for a parameter vector $u \in \mathbb{R}^{n_u \times 1}$,
\begin{align*} 
    \underset{u}{\text{minimize}} & \quad V(u) \\
    \text{subject to} & \quad h(u) = \bm{0}_{n_u \times 1}, \\
    & \quad g(u) \leq \bm{0}_{n_u \times 1},
\end{align*}
a quadratic sub-problem can be formulated as:
\begin{align} 
    \underset{\tilde{u}}{\text{minimize}} & \quad \frac{1}{2} \tilde{u}^T \nabla^2 V(u^l) \tilde{u} + \nabla V(u^l)^T \tilde{u}  \\
    \text{subject to} & \quad \nabla h(u^l)\tilde{u} + h(u^l) = \bm{0}_{n_u \times 1}, \\
    & \quad \nabla g(u^l)\tilde{u} + g(u^l) \leq \bm{0}_{n_u \times 1} \label{eq:linsubproblemIneqConst}, %\bm{0}_{n_u \times 1},
\end{align}
where \( \nabla V(u^l) \) and  $\nabla^2 V(u^l)$ denote the gradient and Hessian of  function \( V \) at \( u^l \), respectively, with  \( u^l \) as the solution estimate at the current time step \( l \) and \( \tilde{u} = u - u^l \). The functions $g: \mathbb{R}^{n_{u} \times 1} \to \mathbb{R}^{n_{u} \times 1}$ and $h: \mathbb{R}^{n_{u} \times 1} \to \mathbb{R}^{n_{u} \times 1}$ denote the inequality and equality constraints, respectively. As in \cite{gonzalez2020nonlinear}, the gradients of the summands of the cost function can be calculated with the matrices $\bm{\Lambda}$ and $\bm{S}$ from Equation~\eqref{eq:nearlyToeplitzMat}
\begin{align*}
    \nabla J_{QP}(\Delta U) &= 2\bm{S}(\bm{\Rho}_k)^T \bm{\mathcal{Q}}_{\text{WT}} \left(\bm{\Lambda}(\bm{\Rho}_k)x_d(k) + \bm{S}(\bm{\Rho}_k) \Delta U(k) - X_{s}\right),\\
    \nabla J_R(\Delta U) &= 2 \bm{\mathcal{R}}_{\text{WT}} \Delta U(k),
    % \nabla J_P(\Delta U_{+}) &= 2\bm{s}(\bm{\Rho}_{k+1}) \bm{P}_{\text{WT}} \left(\bm{\lambda}(\bm{\Rho}_{k+1}) x_d(k) + \bm{s}(\bm{\Rho}_{k+1})  [\Delta U(k)^T \Delta u(k+1)]^T - x_{s}\right),
\end{align*}
and the Hessian as
\begin{align*}
    \nabla^2 J_{QP}(\Delta U) &= 2\bm{S}(\bm{\Rho}_k)^T \bm{\mathcal{Q}}_{\text{WT}} \bm{S}(\bm{\Rho}_k),\\
    \nabla^2 J_R(\Delta U) &= 2 \bm{\mathcal{R}}_{\text{WT}}.
    %\nabla^2 J_P(\Delta U_{+}) &= 2\bm{s}(\bm{\Rho}_{k+1}) \bm{P}_{\text{WT}} \bm{s}(\bm{\Rho}_{k+1}).
\end{align*}
The inequality constraint from Equation~\eqref{eq:linsubproblemIneqConst} is used to express the boundaries of the control inputs. For this velocity based controller the control input limits 
\begin{align*}
    \bm{\Delta u}_{\text{max}} = \begin{bmatrix} \Delta u_{T_g, \text{max}} \quad \Delta u_{\beta_0, \text{max}} \end{bmatrix}^T \in \mathbb{R}^{2 \times 1} \; \text{and}\; 
    \bm{\Delta u}_{\text{min}} = \begin{bmatrix} \Delta u_{T_g, \text{min}} \quad \Delta u_{\beta_0, \text{min}} \end{bmatrix}^T \in \mathbb{R}^{2 \times 1}
\end{align*}
are vectors that contain the respective pitch rate and torque rate limit. All elements of the trajectory $\Delta U$ are limited with
\begin{align*}
    \Delta U \leq \Delta U_{\text{max}} \quad \text{and} \quad \Delta U \geq \Delta U_{\text{min}},
\end{align*}
where $ \Delta U_{\text{max}} = \bm{I}_{n_h-1 \times 1}\otimes \bm{\Delta u}_{\text{max}}$  and $ \Delta U_{\text{min}} = \bm{I}_{n_h-1 \times 1}\otimes \bm{\Delta u}_{\text{min}}$.
%
This can be written using a matrix form as follows:
\begin{align*}
    \begin{bmatrix}
        I_{2 \times n_h} \\
        -I_{2 \times n_h}
    \end{bmatrix} \Delta U \leq
    \begin{bmatrix}
        \Delta U_{\text{max}} \\
        -\Delta U_{\text{min}}\end{bmatrix}
\end{align*}
with the same rate limits applicable to $\Delta \tilde{U}$. This allows reformulating the optimization problem from Equation~\eqref{eq:ObjFunWT} as %
\begin{equation} \label{eq:ObjFunWTQP}
    \begin{aligned} 
        \underset{\Delta \tilde{U}}{\text{minimize }} & \quad 
        \frac{1}{2} \Delta \tilde{U}^T 
        \left(\nabla^2 J_{QP} (\Delta U^l) + \nabla^2 J_R(\Delta U^l) \right) \tilde{U} \\
        & + \left(\nabla J_{QP} (\Delta U^l) + \nabla J_R(\Delta U^l))\right)^T \Delta \tilde{U} \\
        \text{subject to }   &  \begin{bmatrix} I_{2 \times n_h} \\ -I_{2 \times n_h} \end{bmatrix} \Delta U + 
        \begin{bmatrix} - \Delta U_{\text{max}} \\ \Delta U_{\text{min}} \end{bmatrix} \leq 0. 
    \end{aligned}
\end{equation}

The updates of the estimated scheduling, state, and input trajectories are computed using Algorithm~\ref{algo:qLPVMPC}, where \(T_{k,\text{end}}\) denotes the simulation end time. The matrices \(\bm{\Lambda}\) and \(\bm{S}\) are defined as described above, and \(\bm{\Tilde{H}}\) selects the states and inputs to be used for updating the prediction trajectories in the next iteration.
\begin{algorithm}
    \caption{qLMPC}\label{algo:qLPVMPC}
    \begin{algorithmic}[1]
        \Procedure{qLMPC}{$\bm{x}_{s,k}, \bm{x}_{d,k}, \bm{X}_{d,k-1}, \bm{\Rho}_{k}$} 
        \State Initialize plant model, \(\bm{Q}_{\text{WT}}, \bm{R}_{\text{WT}}, \bm{\Rho}_{\text{WT}}\), \(n_h\)
        \State \(k \gets 0\)
        \State \(\bm{\Rho}^0_{0} \gets \mathbf{1}_{n_h} \otimes \bm{\tilde{H}}(\bm{x}_{d,k}, \bm{\Delta u}_{k-1})\)
        \While{$k < T_{k,\text{end}}$} 
        \State \(\bm{X}_{s,k} \gets \mathbf{1}_{n_h} \otimes \bm{x}_{s,k}\)
        \State \(l \gets 0\)
        \While{stop criterion not met}
        \State Compute \(\bm{\Lambda}(\bm{\Rho}_k^l)\) and \(\bm{S}(\bm{\Rho}_k^l)\)
        \State Solve \eqref{eq:ObjFunWTQP} with \(\bm{P}_k^l\) to obtain \(\bm{\Delta U}_k^l\)
        \State Predict \(\bm{X}_k^l = \bm{\Lambda}(\bm{\Rho}_k^l)\bm{x}_{d,k} + \bm{S}(\bm{\Rho}_k^l) \bm{\Delta U}_k^l\)
        \State \(\bm{\Rho}^{l+1}_{k} \gets \bm{\tilde{H}}(\bm{X}_k^l, \bm{\Delta U}_k^l)\)
        \State \(l \gets l + 1\)
        \EndWhile
        \State Apply \(\bm{u}_k\) to the system
        \State \(\bm{\Rho}^0_{k+1} \gets \bm{\tilde{H}}(\bm{X}_k^l, \bm{\Delta U}_k^l)\)
        \State \(k \gets k + 1\)
        \EndWhile
        \State \textbf{return}
        \EndProcedure
    \end{algorithmic}
\end{algorithm}

The controller algorithm requires the following inputs:
\begin{enumerate}
    \item Reference values \(\bm{x}_{s,k}\): the steady-state generator torque and blade pitch, calculated from the measured or estimated wind speed.
    \item Current state \(\bm{x}_{d,k}\).
    \item Previously estimated state trajectory \(\bm{X}_{d,k-1}\).
    \item Scheduling parameter trajectory \(\bm{\Rho}_k\).
\end{enumerate}

Note that the plant model requires turbine parameters used to construct the system matrices \(\bm{A}_d\) and \(\bm{B}_d\), as well as lookup tables for the torque and thrust coefficients \(c_Q\) and \(c_T\). These parameters must be supplied to the controller, which is implemented as a MATLAB function integrated into the Simulink controller model. All required data is stored in a MATLAB data dictionary. In this thesis, the same data is used in both the Simulink controller and plant models. The results for two wind test cases, a stepped wind speed profile and a turbulent wind signal at above-rated conditions, are presented and discussed in the next section.